\documentclass[12pt]{article}
\usepackage[utf8]{inputenc}
\usepackage[T1]{fontenc}
\usepackage{hyperref}
\title{Tepecano Language Resources for K-12 Curriculum Development}
\author{Compiled by Claude Sonnet 5 (claude-sonnet-5) for Andr\'as Kornai}
\date{}
\begin{document}
\maketitle

\section{Introduction}
Tepecano (ISO 639-3 tep) is an extinct Uto-Aztecan (Southern Tepiman) language once spoken at Azqueltan (also spelled Atzqueltlan), Jalisco, Mexico. The last known speaker was Lino de la Rosa, born 1895 and still living in 1980. Closely related Southern Tepehuan is a different language.

\section{Primary Educational Resources}
No school materials or primers were found.

\section{Linguistic and Reference Materials}
\begin{itemize}
\item John Alden Mason's grammar, ``Tepecano, a Piman Language of Western Mexico'', \emph{Annals of the New York Academy of Sciences} 25, pp.\ 309--416. Sources date it to 1916 or 1917. It is the main documentation, and I found no modern grammar. \url{https://grambank.clld.org/sources/sMasonTepecano}
\item Tepecano materials in the John Alden Mason Papers at the American Philosophical Society (fieldwork 1911--1913). \url{https://indigenousguide.amphilsoc.org/entry/39510}
\item Penn Museum finding aid for Mason's Tepecano work. \url{https://findingaids.library.upenn.edu/records/UPENN_MUSEUM_PU-MU.1139}
\item Wikipedia overview. \url{https://en.wikipedia.org/wiki/Tepecano_language}
\end{itemize}

\section{Religious and Community Resources}
None found. Mason also collected native prayers (an article of about 1918), but I did not see a link to it.

\section{Digital and Audio Resources}
None found.

\section{Curriculum Development Context}
Material for a revitalization curriculum is limited to Mason's early twentieth-century grammar and the archival fieldnotes. Any heritage programme would have to be rebuilt from these sources, preferably with the descendant community, and nothing found shows a current programme.
\end{document}